% EVAL FIXTURE — see b1.tex. This file holds the material the Step-1 panel needs and
% will never see: the reason the approach is convincing, the preliminary result, and the
% PI's claim to the problem. A passing review says so and names what to lift into B1 and
% what it displaces from the five pages.
\documentclass[11pt,a4paper]{article}
\usepackage[b2]{ercformatting}
\usepackage{ercreview}

\begin{document}
\maketitle
\startsectiona

\section{Project design}
Three work packages on a shared population-size gradient of 24 populations spanning
40 to 4,000 individuals, so that O1's threshold is estimated rather than assumed.

\section{The result this project starts from}
In a 2025 pilot on six populations we estimated selection gradients and effective size
in the same season for the first time. In the three smallest populations the gradient
on cold tolerance was negative; in the three largest it was positive, as the standard
model predicts. Six populations cannot locate a threshold, but they establish that the
sign change is real and not a census artefact -- which is the fact that makes the
present proposal answerable rather than speculative.

\section{Why this group}
The paired-estimation method is ours: we developed the effective-size estimator in
2022 and the field protocol that makes same-season paired estimation possible in 2024.
The two transects whose disagreement motivates the project are both run by this group,
so the twenty-year census series is available without re-collection.

\startsectionb

\subsection{WP1: Locating the threshold}
\emph{Objective.} A selection gradient and an effective size for each of the 24
populations, in the same season.
\emph{Rationale.} O1 asks whether the gradient reverses below a threshold, and a
threshold cannot be located from six populations.

\paragraph{Task 1.1: Paired estimation across the gradient.}
To locate the threshold we apply the paired protocol to all 24 populations in a single
season. Single-season sampling removes the between-year climate variance that would
otherwise be confounded with population size. This yields the gradient-by-size curve
that O1 is a claim about.

\subsection{Risk assessment}
A single season loses the gradient if that season is climatically atypical. Two seasons
are budgeted, the second conditional on a climate-anomaly test, so an atypical first
season costs a year and not the objective. The risk is worth taking because
multi-season sampling at 24 populations exceeds any single-PI budget.

\startappendix
\begin{fundingtable}
No funding & & & & & \\
\end{fundingtable}

\end{document}
